\section{Regularization, unfolding, and Jacquet adjointness}
\label{jp-section}

We prove the analytic passage from the group operators to the
Jacquet pairing.  The geometric integration formula and the
integrable majorant are the theorems of
Sections~\ref{jm-section} and~\ref{jg-section}.
In particular, none of the estimates in this section assumes
absolute convergence of an unweighted integral over the
unipotent group.

\subsection{The transform and the differentiability class}

Fix finite dimensional complex Hilbert spaces \(W',W\), continuous
unitary representations \(\pi,\pi'\) of \(K_z,K_{z'}\), and
\(E=\operatorname{Hom}_{\mathbb C}(W',W)\).
Use the Hilbert--Schmidt inner product on \(E\), linear in the
first variable.  Replacing its norm by the operator norm changes
the estimates below only by constants depending on the dimensions.

Let \(S_z=N_rK_z\) and \(S_{z'}=N_rK_{z'}\).
Each element has a unique expression \(nk\); the two factors have
trivial intersection because \(N_r\) is unipotent and \(K_z\)
fixes all the outer vectors.  Both groups are unimodular, and
\(dn\,dk\) is Haar measure: compact conjugation preserves \(dn\).
We write \(P=S_z\times S_{z'}\), with product Haar measure \(dp\).
The coordinates \(p=((k_1,n_1),(k_2,n_2))\) are those of
Theorem~\ref{jg-intformula}; their order is retained in formulas.
Set
\[
 \begin{split}
 s_p(g)&=k_1^{-1}n_1^{-1}g n_2k_2,\\
 \rho(p)F(g)&=\pi(k_1)F(s_p(g))\pi'(k_2^{-1}),\\
 \chi_\tau(p)&=\exp\!\bigl(i(\tau_1\lambda_z(n_1)
                              -\tau_2\lambda_{z'}(n_2))\bigr),
                 \qquad \tau\in\mathcal Q=(0,\infty)^2 .
 \end{split}
\]
The characters \(\lambda_z,\lambda_{z'}\) are additive on \(N_r\)
and invariant under their compact stabilizers.  Thus \(\chi_\tau\)
is a unitary character of \(P\).
For the product inherited from \(S_z\times S_{z'}\),
\[
                       \rho(p)\rho(p')=\rho(pp').
\]
This follows directly by multiplying \(n_1k_1\) and \(n'_1k'_1\),
and similarly on the right.  The transformations of \(E\) supplied
by the compact representations preserve its Hilbert--Schmidt norm.
Each \(\rho(p)\) is therefore unitary on \(L^2(G_r,E)\).

For \(X\) in the real Lie algebra of \(G_r\), define
\[
 L_XF(g)=\left.\frac d{dt}F(e^{-tX}g)\right|_{t=0},
 \qquad
 R_XF(g)=\left.\frac d{dt}F(ge^{tX})\right|_{t=0}.
\]
For words \(u=(X_1,\ldots,X_i)\), \(v=(Y_1,\ldots,Y_j)\),
let \(L_u=L_{X_1}\cdots L_{X_i}\) and
\(R_v=R_{Y_1}\cdots R_{Y_j}\), with the empty word equal to
the identity.  Left and right derivatives commute.
Choose the compact unit ball in the Lie algebra consisting of
matrices whose entries have absolute value at most one, and set
\[
 \|F\|_{A,m}
 =\max_{i+j\leq m}\ 
      \sup_{\substack{g\in G_r\\X_a,Y_b\text{ in that ball}}}
        M(g)^A\|L_uR_vF(g)\|.
\]
Write \(\mathcal C_A^m\) for the \(C^m\) functions for which this
quantity is finite.  This is the differentiability and decay class
in the source: a \(C^m\) function on the embedded matrix Lie group
has local \(C^m\) ambient extensions, by extending in local
submanifold coordinates, and conversely such extensions give the
intrinsic derivatives used here.
Throughout the section,
\[
                         A\geq B(k,r).
\]
Theorem~\ref{jg-lone} implies
\(\mathcal C_A^0\subset L^1(G_r,E)\).

Define, wherever the integral is absolutely convergent,
\[
 J_\tau F(g)=\int_P\chi_\tau(p)\rho(p)F(g)\,dp,\qquad
 \Phi_{F,F'}(p)=\int_{G_r}
                  \langle\rho(p)F(g),F'(g)\rangle\,dg.
\]
Let \(X(p)=\Xi(n_1k_1)\Xi(n_2k_2)\), for the function
\(\Xi\) constructed in Section~\ref{jm-section}.
Its two proved properties give
\begin{align}
 \int_{G_r}\|\rho(p)F(g)\|\|F'(g)\|\,dg
       &\leq C\|F\|_{A,0}\|F'\|_{A,0}X(p),
                                      \label{jp-coef}\\
 \int_P\frac{X(p)}
       {(1+|\lambda_z(n_1)|)^2(1+|\lambda_{z'}(n_2)|)^2}\,dp
       &<\infty .                     \label{jp-xi}
\end{align}
In particular \(\Phi\) is an ordinary absolutely convergent
group integral for every \(p\).

\begin{lemma}\label{jp-orbital}
For \(F\in\mathcal C_A^0\), the integral \(J_\tau F(g_q)\)
is absolutely convergent for almost every \(q\) with respect to
\(\nu\), on one common set valid for every \(\tau\in\mathcal Q\).
On that set it is continuous in \(\tau\).
It has a jointly measurable version in \((\tau,q)\).
\end{lemma}

\begin{proof}
The integration formula and unitarity on \(E\) give
\[
 \int_Q\int_P\|\rho(p)F(g_q)\|\,dp\,d\nu(q)
                         =\int_{G_r}\|F(g)\|\,dg<\infty .
\]
Tonelli supplies the common full measure set.
The integrand's norm is independent of \(\tau\), so dominated
convergence gives continuity there.
The map \((\tau,q,p)\mapsto\chi_\tau(p)\rho(p)F(g_q)\) is
measurable.  Integration of a measurable finite dimensional
vector valued function, with value zero assigned where its
norm integral is infinite, is measurable; this follows
componentwise from truncation and the scalar integration
theorem.  This proves the final assertion.
\end{proof}

We use this version to define
\[
 N_\tau(F)=\int_Q\|J_\tau F(g_q)\|^2\,d\nu(q)\in[0,\infty],
 \qquad
 \beta_\tau(F,F')=
       \int_Q\langle J_\tau F(g_q),J_\tau F'(g_q)\rangle\,d\nu(q)
\]
when both squared norms are finite.
The eventual pairing in the source is \(\beta=\beta_{(1,1)}\).

\subsection{Dilation estimates and the regularized form}

For \(t=(t_1,t_2)\in\mathcal Q\), put
\[
 D_tF(g)=F(a_{t_1}^{-1}g a_{t_2}),\qquad
 U_tF(g)=F(a_{t_1}g a_{t_2}^{-1}).
\]
Here \(U_t=D_{(t_1^{-1},t_2^{-1})}\).

\begin{lemma}\label{jp-dilation-seminorm}
On every compact rectangle in \(\mathcal Q\),
\(\|D_tF\|_{A,m}\leq C\|F\|_{A,m}\).
On \([1/2,2]^2\) one also has
\begin{equation}\label{jp-dilation-cont}
 \|D_tF-F\|_{A,4}
       \leq C(|t_1-1|+|t_2-1|)\|F\|_{A,5}.
\end{equation}
For \(|\alpha|\leq m\), the ordinary parameter derivatives
\(\partial_t^\alpha U_tF\) have decay norm of order zero at
most \(C_\alpha\|F\|_{A,m}\), uniformly on a compact rectangle.
\end{lemma}

\begin{proof}
The Frobenius norm is submultiplicative, and the matrices
\(a_t,a_t^{-1}\) are uniformly bounded on such a rectangle.
Consequently \(M(a_{t_1}^{-1}g a_{t_2})\) and \(M(g)\) are
comparable there, uniformly in \(g\).
The chain rule gives
\[
 L_XD_tF=D_t L_{\operatorname{Ad}(a_{t_1}^{-1})X}F,\qquad
 R_YD_tF=D_t R_{\operatorname{Ad}(a_{t_2}^{-1})Y}F.
\]
Iterating these identities, and using multilinearity in the
Lie algebra elements, proves the first bound: the indicated
adjoint maps have uniformly bounded coefficients.

Let \(H\) be the diagonal matrix with entries
\(r,r-1,\ldots,1,0,\ldots,0,-1,\ldots,-r\).
Then \(a_t=e^{(\log t)H}\), \(H\) belongs to the real Lie
algebra, and
\[
       \partial_{t_1}D_tF=t_1^{-1}D_tL_HF,\qquad
       \partial_{t_2}D_tF=t_2^{-1}D_tR_HF.
\]
The corresponding formulas for \(U_t\) have the opposite
signs and the same bounded scalar factors.
Higher parameter derivatives are finite sums of translated
Lie derivatives, with smooth bounded coefficients.
This proves their asserted decay estimates.

Apply one parameter derivative also to the displayed formulas
for \(L_uR_vD_tF\), with \(|u|+|v|\leq4\).
It differentiates either an adjoint coefficient or the
translated function.  The latter term uses at most five
Lie derivatives of \(F\).  All terms are bounded, after
multiplication by \(M(g)^A\), by \(C\|F\|_{A,5}\).
Integrate first along \(t_1\) and then along \(t_2\) from
\((1,1)\) to \(t\).  Taking the suprema defining the
seminorm gives \eqref{jp-dilation-cont}.
\end{proof}

Choose a real smooth function \(\chi\), supported in
\((1/2,2)\), with \(\int_{\mathbb R}\chi(t)\,dt=1\).
It need not be nonnegative.  Write
\[
 \delta_t p=((k_1,a_{t_1}n_1a_{t_1}^{-1}),
             (k_2,a_{t_2}n_2a_{t_2}^{-1})),\qquad
 \lambda_1=\lambda_z(n_1),\quad\lambda_2=\lambda_{z'}(n_2).
\]
Define the compactly supported \(C^4\) function
\[
 a_{F,F',p}(t)=
       \chi(t_1)\chi(t_2)(t_1t_2)^{\kappa_r}
                                  \Phi_{F,F'}(\delta_t p),
\]
extended by zero outside \(\mathcal Q\).
For \(F,F'\in\mathcal C_A^4\), define
\begin{equation}\label{jp-reg-definition}
 B_\tau(F,F')
 =\int_P\frac{1}{(1+\tau_1^2\lambda_1^2)
                    (1+\tau_2^2\lambda_2^2)}
    \int_{\mathbb R^2}
      (1-\partial_{t_1}^2)(1-\partial_{t_2}^2)
          a_{F,F',p}(t)\,
      e^{i(\tau_1t_1\lambda_1-\tau_2t_2\lambda_2)}\,dt\,dp .
\end{equation}
This is exactly the source's regularized pairing.

\begin{proposition}\label{jp-reg-bound}
The integral \eqref{jp-reg-definition} is absolutely convergent.
It is sesquilinear in \(F,F'\), continuous in \(\tau\in\mathcal Q\),
and satisfies, locally uniformly in \(\tau\),
\begin{equation}\label{jp-Bbound}
       |B_\tau(F,F')|\leq C_\tau\|F\|_{A,4}\|F'\|_{A,4}.
\end{equation}
\end{proposition}

\begin{proof}
The substitution \(g=a_{t_1}h a_{t_2}^{-1}\) in the group
integral gives
\begin{equation}\label{jp-phase-translate}
 \Phi_{F,F'}(\delta_t p)
                  =\Phi_{U_tF,U_tF'}(p).
\end{equation}
It uses both unimodularity of \(G_r\) and commutation of
\(a_t\) with the compact stabilizers.
For fixed \(p\), the derivatives in \(t\) on the right may
be taken under the group integral.  Indeed every derivative
of total order at most four is bounded, on the fixed compact
rectangle, by a constant times
\[
                 M(s_p(g))^{-A}M(g)^{-A}.
\]
This is integrable by \eqref{jp-coef}.
The same bound applies to the continuous highest derivatives,
so repeated dominated convergence proves the asserted \(C^4\)
regularity, including the mixed derivative of order \((2,2)\).

After differentiation, apply \eqref{jp-coef} to each of the
translated derivatives, and then the preceding lemma.
For \(0\leq i,j\leq2\) this yields
\[
 |\partial_{t_1}^i\partial_{t_2}^j a_{F,F',p}(t)|
 \leq C\|F\|_{A,4}\|F'\|_{A,4}X(p)
                      {\bf1}_{[1/2,2]^2}(t).
\]
On a compact subset of \(\mathcal Q\),
\[
 (1+\tau_j^2\lambda_j^2)^{-1}
                         \leq C(1+|\lambda_j|)^{-2}.
\]
Equation~\eqref{jp-xi} proves absolute convergence and the
bound.  It also supplies a common dominating function for
continuity in \(\tau\).  Sesquilinearity follows from that
of the integrands, with all the integrals now justified.
\end{proof}

\subsection{Averaging in the two characters}

For a real \(C^2\) function \(w\) of compact support in
\((0,\infty)\), write
\[
                \widehat w(x)=\int_{\mathbb R}w(\tau)e^{i\tau x}\,d\tau.
\]
Twice integrating by parts gives
\[
          |\widehat w(x)|\leq C_w(1+|x|)^{-2}.
\]
For example combine the bounds \(\|w\|_1\) and
\(|x|^{-2}\|w''\|_1\); the boundary terms vanish by compact
support.
For two such functions set
\(\Omega(p)=\widehat w_1(\lambda_1)\widehat w_2(-\lambda_2)\).

\begin{proposition}\label{jp-reg-average}
For \(F,F'\in\mathcal C_A^4\),
\begin{equation}\label{jp-average}
 \int_{\mathcal Q}w_1(\tau_1)w_2(\tau_2)B_\tau(F,F')\,d\tau
                           =\int_P\Omega(p)\Phi_{F,F'}(p)\,dp .
\end{equation}
The integral on the right is absolutely convergent.
\end{proposition}

\begin{proof}
Its absolute convergence follows from
\eqref{jp-coef}, \eqref{jp-xi}, and Fourier decay.
For each fixed \(p\), integration by parts in \(t_1,t_2\)
gives
\[
 \begin{split}
 &\frac{1}{(1+\tau_1^2\lambda_1^2)(1+\tau_2^2\lambda_2^2)}
    \int (1-\partial_{t_1}^2)(1-\partial_{t_2}^2)a_p(t)
                         e^{i(\tau_1t_1\lambda_1-\tau_2t_2\lambda_2)}\,dt\\
 &\hspace{25mm}
       =\int a_p(t)e^{i(\tau_1t_1\lambda_1-\tau_2t_2\lambda_2)}\,dt .
 \end{split}
\]
The right side is asserted here for fixed \(p\); no unweighted
interchange of \(p\) and \(t\) is being made.
The absolute bound used in Proposition~\ref{jp-reg-bound}
allows averaging the left side in \(\tau\) under the \(p\)
integral.  For fixed \(p\), all \(t,\tau\) integrals are
over compact sets, so their interchange gives
\[
 \int_P\int_{\mathbb R^2}
       \chi(t_1)\chi(t_2)(t_1t_2)^{\kappa_r}
       \Phi_{F,F'}(\delta_t p)
       \widehat w_1(t_1\lambda_1)
       \widehat w_2(-t_2\lambda_2)\,dt\,dp .
\]
This new double integral is absolutely convergent.
To see this, for each \(t\) in the support change variables
\(p'=\delta_t p\).
Proposition~\ref{jg-dilation} gives
\(dp=(t_1t_2)^{-\kappa_r}dp'\), and its character identity
replaces \(t_j\lambda_j\) by the corresponding \(\lambda_j'\).
The absolute integral is therefore bounded by
\[
 \|\chi\|_1^2\int_P|\Omega(p')\Phi_{F,F'}(p')|\,dp'<\infty .
\]
Fubini now applies.  The same substitution without absolute
values and \((\int\chi)^2=1\) prove \eqref{jp-average}.
\end{proof}

\begin{proposition}\label{jp-unfold}
For \(F\in\mathcal C_A^4\) and nonnegative \(w_1,w_2\) as above,
\begin{equation}\label{jp-unfold-identity}
 \int_{\mathcal Q}w_1(\tau_1)w_2(\tau_2)N_\tau(F)\,d\tau
                  =\int_P\Omega(p)\Phi_{F,F}(p)\,dp<\infty .
\end{equation}
The equality regards the nonnegative real number on the left
as a complex number on the right.
\end{proposition}

\begin{proof}
Fix \(q\) in the common full measure set of
Lemma~\ref{jp-orbital} and put \(f_q(p)=\rho(p)F(g_q)\).
This is an \(E\)-valued \(L^1(P)\) function.
Expanding the squared norm of its character integral gives
an absolutely convergent double integral.
In that integral set the second variable equal to \(p\)
and the first to \(pp'\).
Left invariance of Haar measure, the character identity,
and unitarity on \(E\) give
\begin{equation}\label{jp-fixed-q}
 \|J_\tau F(g_q)\|^2
 =\int_P\int_P \chi_\tau(p')
     \langle\rho(p')F(s_p(g_q)),F(s_p(g_q))\rangle\,dp'\,dp .
\end{equation}
The compact factors attached to \(p\) act unitarily.
They transform the two entries of the inner product on
the right into \(\rho(pp')F(g_q)\) and \(\rho(p)F(g_q)\).
The absolute double integral before inserting the character
is at most \(\|f_q\|_1^2\).
Thus averaging \eqref{jp-fixed-q} in \(\tau\) gives the same
formula with \(\chi_\tau(p')\) replaced by \(\Omega(p')\).

It remains to justify integration in \(q\); the preceding
bound need not be integrable in \(q\).
Instead use the Fourier decay already obtained.
The integration formula for \(G_r\), followed by Tonelli, gives
\[
 \begin{split}
 &\int_Q\int_P\int_P|\Omega(p')|
       \|\rho(p')F(s_p(g_q))\|\,\|F(s_p(g_q))\|
                        \,dp'\,dp\,d\nu(q)\\
 &\quad=\int_P|\Omega(p')|
               \int_{G_r}\|\rho(p')F(g)\|\,\|F(g)\|\,dg\,dp'
 <\infty ,
 \end{split}
\]
by \eqref{jp-coef} and \eqref{jp-xi}.
This proves that the averaged right side of
\eqref{jp-fixed-q} is integrable in \(q\).
On the left, \(w_1w_2\geq0\), so Tonelli applies even before
finiteness is known.  Fubini on the absolutely integrable
right side now proves \eqref{jp-unfold-identity}.
\end{proof}

\subsection{Equality at every positive character}

\begin{theorem}\label{jp-pairing-equality}
For every \(F,F'\in\mathcal C_A^5\) and every
\(\tau\in\mathcal Q\), the two squared norms
\(N_\tau(F),N_\tau(F')\) are finite and
\begin{equation}\label{jp-beta-equals}
                         \beta_\tau(F,F')=B_\tau(F,F').
\end{equation}
In particular
\[
       |\beta(F,F')|\leq C\|F\|_{A,4}\|F'\|_{A,4}.
\]
\end{theorem}

\begin{proof}
First take only \(F\in\mathcal C_A^4\), and put
\(b(\tau)=B_\tau(F,F)\).
The two averaging propositions say
\[
       \int w_1w_2N_\tau(F)\,d\tau=\int w_1w_2b(\tau)\,d\tau
\]
for every pair of nonnegative test functions.
The left side is finite.
On any compact rectangle in \(\mathcal Q\), choose test
functions positive on that rectangle; this shows that
\(\tau\mapsto N_\tau(F)\) is locally integrable and finite
almost everywhere.
At its Lebesgue points, use products of nonnegative smooth
approximate identities in the displayed equality.
Their support can be kept inside \(\mathcal Q\).
The Lebesgue differentiation theorem on the left and continuity
of \(b\) on the right give
\[
                           N_\tau(F)=b(\tau)
                           \quad\hbox{for almost every }\tau .
\]
In particular \(b\) is real and nonnegative almost everywhere,
and hence everywhere by continuity.

For any fixed \(\tau_0\), choose a sequence of such equality
points tending to \(\tau_0\).
The pointwise continuity in Lemma~\ref{jp-orbital} and
Fatou's lemma give
\begin{equation}\label{jp-diagonal-domination}
                    N_{\tau_0}(F)\leq b(\tau_0).
\end{equation}
Thus the norm is finite at every positive character even
for \(F\in\mathcal C_A^4\).  Moreover, locally uniformly in
\(\tau\),
\[
              \|J_\tau F(g_\cdot)\|_{L^2(Q,\nu)}
                                      \leq C\|F\|_{A,4}.
\]
Linearity holds because the defining orbital integrals are
absolutely convergent on common full measure sets.

We next use the fifth derivative.
Changes of variables \(p=\delta_t p'\), and then
\(q'=S_{t_1t_2}q\), give
\begin{align}
 J_\tau(D_tF)(g)
     &=(t_1t_2)^{\kappa_r}
           J_{(\tau_1t_1,\tau_2t_2)}F(a_{t_1}^{-1}g a_{t_2}),
                                                   \label{jp-J-dilation}\\
 N_\tau(D_tF)
     &=(t_1t_2)^{\kappa_r}
                           N_{(\tau_1t_1,\tau_2t_2)}(F).
                                                   \label{jp-N-dilation}
\end{align}
These equalities use precisely the two pushforward exponents
in Proposition~\ref{jg-dilation}.
They are valid on common full measure sets in \(q\);
the radial scaling preserves null sets.

For \(F\in\mathcal C_A^5\), the estimate just proved and
\eqref{jp-dilation-cont} show
\[
 \|J_\tau(D_tF)(g_\cdot)-J_\tau F(g_\cdot)\|_{L^2}
 \leq C\|D_tF-F\|_{A,4}
                 \longrightarrow0\qquad(t\longrightarrow(1,1)).
\]
Consequently \(N_\tau(D_tF)\to N_\tau(F)\).
Equation~\eqref{jp-N-dilation} now proves continuity of
\(\tau'\mapsto N_{\tau'}(F)\) at \(\tau\).
Since \(\tau\) was arbitrary, it is continuous on
\(\mathcal Q\), and its almost everywhere equality with
\(b\) is therefore equality everywhere.

Finally both sides of \eqref{jp-beta-equals} are
sesquilinear on \(\mathcal C_A^5\).
The integral defining the left side is absolutely convergent
by Cauchy--Schwarz and the finite diagonal norms.
Equality of the diagonal values, followed by complex
polarization, proves \eqref{jp-beta-equals}.
The last bound is \eqref{jp-Bbound} at \((1,1)\).
\end{proof}

\subsection{Adjoints without differentiation of an orbital integral}

Let \(U\subset G_r\) be open and stable under left \(S_z\)
and right \(S_{z'}\).
Consider the finite order differential operator
\[
              \Theta F=\sum_T A_T(g)R_{w_T}F,
 \qquad
 \Theta^\dagger F'
       =\sum_T(-1)^{|w_T|}
                     R_{\operatorname{rev}(w_T)}(A_T(g)^*F').
\]
The coefficients and their pointwise adjoints are smooth
on \(U\), and \(m=\max_T|w_T|\).
Assume that \(\Theta\rho(p)F=\rho(p)\Theta F\) on \(U\)
for \(C^m\) functions \(F\).
The support of a function is its closed support in \(G_r\).

\begin{lemma}\label{jp-ibp}
If a \(C^1\) scalar function \(h\) and \(R_Xh\) are in
\(L^1(G_r)\), then \(\int R_Xh\,dg=0\).
Consequently integration by parts against a real Lie
derivative has no boundary term whenever the three products
\(\langle f,f'\rangle\),
\(\langle R_Xf,f'\rangle\), and
\(\langle f,R_Xf'\rangle\) are integrable.
\end{lemma}

\begin{proof}
The fundamental theorem of calculus gives, for \(t>0\),
\[
       h(ge^{tX})-h(g)=\int_0^t(R_Xh)(ge^{sX})\,ds.
\]
The absolute double integral on the right is
\(t\|R_Xh\|_1\), by right invariance of Haar measure.
Fubini therefore applies.  Integration of the left side
gives zero, and the right side gives \(t\int R_Xh\,dg\).
For the second assertion use the product rule for the
inner product, which is differentiated in a real parameter.
\end{proof}

\begin{theorem}\label{jp-adjoint}
Suppose \(F,F'\) are \(C^{m+5}\), supported in \(U\), and
\[
       F,F',\Theta F,\Theta^\dagger F'\in\mathcal C_A^5 .
\]
For every \(p\in P\), every term \(T\), and
\(0\leq i\leq j\leq|w_T|\), assume the actual integrability
condition
\begin{equation}\label{jp-intclause}
 \int_{G_r}
   \bigl\|R_{\operatorname{rev}(w_T[0:i])}(A_T^*F')(g)\bigr\|
   \bigl\|R_{w_T[j:]}(\rho(p)F)(g)\bigr\|\,dg<\infty .
\end{equation}
Then, for every \(\tau\in\mathcal Q\),
\begin{equation}\label{jp-adjoint-identity}
            \beta_\tau(\Theta F,F')
                         =\beta_\tau(F,\Theta^\dagger F').
\end{equation}
\end{theorem}

\begin{proof}
The product \(A_T^*F'\), extended by zero outside \(U\),
is \(C^{m+5}\) on all of \(G_r\).
At every point outside \(U\), the closed support hypothesis
gives a neighborhood on which \(F'\) is zero, so no
regularity of \(A_T\) outside \(U\) is needed.

Fix \(p,T\), and transfer the letters of \(w_T\) one at a
time from \(\rho(p)F\) to \(A_T^*F'\).
At the \(j\)-th transfer, the three product norms needed in
Lemma~\ref{jp-ibp} are exactly the cases
\((i,j)=(j,j+1),(j,j),(j+1,j+1)\) of
\eqref{jp-intclause}.  Thus each transfer is legitimate,
and
\[
 \int\langle R_{w_T}\rho(p)F,A_T^*F'\rangle\,dg
   =(-1)^{|w_T|}
      \int\langle\rho(p)F,
                 R_{\operatorname{rev}(w_T)}(A_T^*F')\rangle\,dg .
\]
The reversal of the word and the sign are both necessary.
Summing and using commutation on \(U\), while \(F'\) is
zero outside \(U\), proves
\[
                  \Phi_{\Theta F,F'}(p)
                               =\Phi_{F,\Theta^\dagger F'}(p).
\]
The same identity holds at \(\delta_t p\).
It therefore holds for their regularized brackets and their
parameter derivatives, and gives equality of the two
\(B_\tau\) forms.  Theorem~\ref{jp-pairing-equality} now
gives \eqref{jp-adjoint-identity}.
No derivative of the unweighted orbital integral was used.
\end{proof}

\subsection{The polynomial coefficient criterion}

To check \eqref{jp-intclause} in the setting used by the
source, put \(v_a=e_a\) for \(a<r\), \(v_r=z\), and use
\(v'_r=z'\) on the right.
Let
\[
 D_l(g)=\det\bigl(h_r(gv'_b,v_a)\bigr)_{0\leq a,b<l},
 \quad D_0=1,\qquad
 \mathcal G(g)=|D_{r+1}(g)|^2-TT'|D_r(g)|^2 .
\]
Here this notation for \(D_l\) refers to these Bessel
minors, not the flag minors of Section~\ref{jm-section}.
The function \(\mathcal G\) is a real polynomial in the
real and imaginary matrix entries.
Every \(S_z,S_{z'}\) change of the indicated frames is
triangular with diagonal one.  Hence every \(D_l\),
and therefore \(\mathcal G\), is invariant under those
left and right actions.
It follows that \(U=\{\mathcal G>0\}\) is stable and open.

\begin{corollary}\label{jp-poly-adjoint}
Let \(d\geq0\), \(\epsilon>0\), and suppose that for every
coefficient \(A_T,A_T^*\), every left and right word of
total length at most \(m+5\), and every
\(\mathcal G(g)\geq\epsilon\), its derivative has norm at
most \(C M(g)^d\), uniformly for Lie algebra letters in
the fixed unit ball.
Assume that \(\Theta\) is smooth and commutes with \(\rho\)
on \(\{\mathcal G>0\}\).
If \(F,F'\in\mathcal C_{A+d}^{m+5}\) vanish whenever
\(\mathcal G<\epsilon\), then
\[
                         \beta(\Theta F,F')
                                      =\beta(F,\Theta^\dagger F').
\]
\end{corollary}

\begin{proof}
Continuity of \(\mathcal G\) places both closed supports
in \(\{\mathcal G\geq\epsilon\}\subset U\).
The Leibniz rule gives
\[
                  F,F',\Theta F,\Theta^\dagger F'
                                           \in\mathcal C_A^5 .
\]
Indeed each differentiated term is bounded by
\(C M^d M^{-(A+d)}\), and uses at most \(m+5\)
derivatives of either factor.  Fixed letters occurring in
\(\Theta\) are handled by scaling them into the unit ball.
All derivatives vanish locally where the function is
outside its closed support.

The first factor in \eqref{jp-intclause} is consequently
bounded by \(C M(g)^{-A}\).
For fixed \(p\), derivatives of the second factor are
finite linear combinations of derivatives of \(F\) at
\(s_p(g)\), with fixed conjugated Lie algebra letters.
Submultiplicativity of \(M\) gives
\(M(s_p(g))^{-(A+d)}\leq C_pM(g)^{-(A+d)}\).
Their product is therefore at most
\(C_p M(g)^{-2A-d}\leq C_pM(g)^{-A}\).
This is integrable by Theorem~\ref{jg-lone}, proving
every case of \eqref{jp-intclause}.
The preceding theorem applies.
\end{proof}

In particular, polynomial matrix coefficients satisfy this
criterion: invariant differentiation acts by linear vector
fields on matrix entries and does not increase polynomial
degree.  The criterion also covers coefficients
\(P(g)\mathcal G(g)^{-q}\) on \(U\), for a polynomial \(P\)
and integer \(q\geq0\).  After \(j\) derivatives, the
Leibniz and chain rules give denominators
\(\mathcal G^{q+i}\), \(i\leq j\), and numerator degree at
most \(\deg P+j\deg\mathcal G\).
On \(\mathcal G\geq\epsilon\) the denominators are bounded
away from zero, giving the required polynomial growth for
all the finitely many derivatives.

This proves the source's pairing bound, equality with the
regularized pairing, adjointness with its full integrability
clause, and its polynomial coefficient version.
The fifth derivative has a specific role: it gives continuity
of dilations in the norm using four derivatives and upgrades
the almost everywhere identity in the characters to an
identity at \((1,1)\).
Absolute convergence of \(J_\tau F(g)\) at every group
element has neither been assumed nor concluded.
All source correspondences are recorded in the accompanying
trace and manifest \cite{alpo-uniform-source}.


